\chapter{语法速查}
\label{app:quickref}

构造按阅读顺序分组：先是声明，然后是迭代、分块选择与数据搬运，再是语句、协调与诊断。第 3 章说明每个构造的含义及其设计依据；本附录的用途是查写法，无需重读示例。

\section{声明与形状}
\begin{table}[htbp]\centering\small
\begin{tabularx}{\linewidth}{@{}lX@{}}
\toprule
构造 & 含义 \\
\midrule
\co@__co__ T [shape] f(T [shape] a)@ & 带形状参数的 Croqtile 函数。 \\
\co@__co__ void@ / \co@__co__ auto@ & 就地结果 / 推导返回类型。 \\
\co@T [a.span] b;@ & 复制形状的缓冲区。 \\
\co@a.span@ / \co@a.span(i)@ & 整体形状 / 单个维度。 \\
\co@x : [7, 8]@ & 具名 \co@mdspan@ 声明。 \\
\co@a = {1, 2, 3}@ & \co@ituple@ 声明。 \\
\co@tensor.at(i, j)@ & 位置访问。 \\
\bottomrule
\end{tabularx}
\end{table}

\section{迭代}
\begin{table}[htbp]\centering\small
\begin{tabularx}{\linewidth}{@{}lX@{}}
\toprule
构造 & 含义 \\
\midrule
\co@parallel p by N@ / \co@parallel p, q by [M, N]@ & 相互独立的实例。 \\
\co@: block@ / \co@: thread@ / \co@: group@ & 并行轴上的标注。 \\
\co@foreach idx in [N]@ / \co@foreach t(1:)@ & 有序迭代 / 切片迭代。 \\
\co@with idx in [a, b]@ & 有作用域的索引绑定。 \\
\co@p # m@ / \co@#p@ & 组合索引 / 轴范围。 \\
\co@cdiv(a, b)@ & 向上取整除法。 \\
\bottomrule
\end{tabularx}
\end{table}

\section{分块选择与数据搬运}
\begin{table}[htbp]\centering\small
\begin{tabularx}{\linewidth}{@{}lXX@{}}
\toprule
构造 & 含义 \\
\midrule
\co@f = dma.copy src => shared@ & 分块搬运语句。 &  \\
\textless{}ROW\textgreater{} & \co@a.chunkat(t)@ & 规则划分中的一个分块。 \\
\co@a.subspan(M, K).at(i, j)@ & 显式范围与分块坐标。 &  \\
\co@a.subspan(M, K).step(sM, sK).at(i, j)@ & 跨步分块选择。 &  \\
\co@a.view(M, N).from(r, c)@ & 带显式元素偏移的窗口。 &  \\
\co@a.span_as([dims])@ & 不搬运数据的形状重解释。 &  \\
\co@f.data@ / \co@f.span@ / \co@dma.any@ & 已搬运分块 / 其形状 / 占位句柄。 &  \\
\co@.async@ / \co@wait f@ / \co@swap@ / \co@rotate@ & 异步传输与句柄轮转。 &  \\
\co@global@ / \co@shared@ / \co@local@ & 存储说明符。 &  \\
\bottomrule
\end{tabularx}
\end{table}

\section{语句修饰符与收缩}
\begin{table}[htbp]\centering\small
\begin{tabularx}{\linewidth}{@{}lX@{}}
\toprule
构造 & 含义 \\
\midrule
\co@dma.transp<...>@ / \co@dma.pad<...>@ & 搬运时变换。 \\
\co@.zfill@ / \co@.swiz<N>@ & 填充与布局修饰符。 \\
\co@mc = mma.fill 0.0@ & 累加器清零。 \\
\co@ma = mma.load src.chunkat(...)@ & 装入操作数。 \\
\co@mma.row.row mc, ma, mb@ & 带布局约定的收缩。 \\
\co@mma.store mc, dst@ & 写回累加器。 \\
\co@.sp@ / \co@.scale@ / \co@mma.scale@ & 收缩变体与阶段栅栏。 \\
\bottomrule
\end{tabularx}
\end{table}

\section{协调、互操作与诊断}
\begin{table}[htbp]\centering\small
\begin{tabularx}{\linewidth}{@{}lX@{}}
\toprule
构造 & 含义 \\
\midrule
\co@if (expr) ... @ & 运行时谓词。 \\
\co@inthreads.async (cond)@ / \co@inthreads (cond)@ & 并发区域 / 顺序过滤。 \\
\co@shared event e[N]@ / \co@trigger e@ / \co@wait e@ & 同步令牌。 \\
\co@parallel.async ...@ / \co@stream s@ / \co@yield@ & 主机调度与提前返回。 \\
\co@__device__ fn()@ / \co@call fn(args)@ & 互操作的 C++ 函数及其调用。 \\
\co@__cpp__("...")@ / \co@__cpp__(R"(...)")@ & 注入的源文本。 \\
\co@#define@ / \co@#if@ / \co@#ifdef@ / \co@#error@ / \co@-D@ & 预处理器。 \\
\co@print!@ / \co@println!@ / \co@print@ / \co@println@ & 编译时输出与运行时输出。 \\
\co@assert@ / \co@croq_assert@ & 设备侧与主机侧断言。 \\
\bottomrule
\end{tabularx}
\end{table}

若要了解构造背后的含义，请回到引入它的小节：初次阅读见第 2 章，设计依据见第 3 章。若要在完整程序中观察某个构造，可从 3.4.1、3.6.1 与 3.8.5 节的示例核函数开始。

